% !TeX document-id = {4a66a092-a797-425f-b588-c0b0f6bcb4d3}
% !TeX TXS-program:compile = txs:///lualatex/ | txs:///biber | txs:///lualatex | txs:///lualatex | txs:///view
% arara: clean: { extensions: ['aux', 'bbl', 'bcf', 'blg', 'log', 'pdf', 'run.xml'] }
% arara: lualatex
% arara: biber
% arara: lualatex
% arara: lualatex
% arara: clean: { extensions: ['aux', 'bbl', 'bcf', 'blg', 'log', 'run.xml'] }
\documentclass[
	DIV=calc,
	final,
	fontsize=11pt,
	paper=landscape
]{scrartcl}
\usepackage[left=1.2cm,right=1.2cm,top=1.8cm,bottom=1.8cm]{geometry}

% Font loading and micro-typography (LuaLaTeX)
\usepackage{fontspec}
\setmainfont{Latin Modern Roman}
\setsansfont{Latin Modern Sans}
\renewcommand{\familydefault}{\sfdefault} % use sans as the default family
\usepackage{microtype}

\usepackage[spanish]{babel}
\usepackage{csquotes}

\usepackage[svgnames]{xcolor}
\usepackage{colortbl}
\usepackage{booktabs}
\usepackage{longtable}
\usepackage{multirow}
\usepackage{array,tabularx}
% ragged-right paragraph column type
\newcolumntype{P}[1]{>{\raggedright\arraybackslash}p{#1}}

\usepackage[
	colorlinks=true,
	urlcolor=Primary,
	linkcolor=Primary,
	anchorcolor=Primary,
	citecolor=Primary
]{hyperref}

\definecolor{Primary}{named}{DarkSlateBlue}
\definecolor{Secondary}{named}{CornflowerBlue}
\definecolor{TeoriaColor}{RGB}{230,240,250}
\definecolor{LabColor}{RGB}{232,245,233}
\definecolor{EvalColor}{RGB}{255,243,224}

\newcommand{\cellteoria}[1]{\cellcolor{TeoriaColor}\parbox[c]{7cm}{\small #1}}
\newcommand{\celllab}[1]{\cellcolor{LabColor}\parbox[c]{7cm}{\small #1}}
\newcommand{\cellevaluacion}[1]{\cellcolor{EvalColor}\parbox[c]{7cm}{\centering\small #1}}
\newcommand{\cellvacio}[1]{\cellcolor{white}\parbox[c]{7cm}{\small #1}}

\usepackage[citestyle=numeric,backend=biber,style=ieee]{biblatex}
\addbibresource{references.bib}

\pagestyle{empty}

\newcommand{\CourseTitle}{Sílabo dosificado de Métodos Numéricos para Ecuaciones Diferenciales Parciales~CM 032 A}
\newcommand{\Professor}{Fidel Jara Huanca}
\newcommand{\Term}{2026-1}
\newcommand{\CourseURL}{https://cm032.github.io}

% PDF metadata and hyperref setup
\hypersetup{
	pdftitle={\CourseTitle},
	pdfauthor={\Professor},
	pdfsubject={Syllabus},
	unicode=true
}

% KOMA-Script heading fonts
% \setkomafont{section}{\normalfont\Large\bfseries}
% \setkomafont{subsection}{\normalfont\large\bfseries}

\begin{document}


\begin{center}
	\Large\bfseries\color{Primary}
	\CourseTitle

	\vspace{0.1cm}

	\small\color{Primary}
	Profesor: \Professor.\qquad
	Ciclo: \Term.\qquad
	Sitio web: \url{\CourseURL}

	\textcolor{Secondary}{\rule{\textwidth}{0.3pt}}
\end{center}

\begin{center}
	\small
	\colorbox{TeoriaColor}{\strut Teoría: (16:00-18:00 horas)}\qquad
	\colorbox{LabColor}{\strut Laboratorio: (17:00-19:00 horas)}\qquad
	\colorbox{EvalColor}{\strut Evaluación}
\end{center}

\begin{longtable}{|c|P{7.2cm}|P{7.2cm}|P{7.2cm}|}

	\hline
	\rowcolor{Primary}
	\multicolumn{1}{|c|}{\textcolor{white}{\bfseries Semana}}                                     &
	\multicolumn{1}{c|}{\textcolor{white}{\bfseries Lunes}}                                       &
	\multicolumn{1}{c|}{\textcolor{white}{\bfseries Miércoles}}                                   &
	\multicolumn{1}{c|}{\textcolor{white}{\bfseries Viernes}}                                       \\
	\hline
	\endhead

	\hline
	\endlastfoot
	1                                                                                             &
	\cellteoria{
		Entrega del sílabo.
		Motivación al curso.
		Prueba de entrada.
	}                                                                                             &
	\celllab{
		Introducción a Python~\cite{python314},
		NumPy~\cite{Harris2020}, SciPy~\cite{Virtanen2020},
		SymPy~\cite{Meurer2017} y Matplotlib~\cite{Hunter2007}.
		Método de diferencia finita para un problema de valor de frontera
		de $2$ puntos~\cite{Köbis2025}.
	}                                                                                             &
	\cellteoria{Aproximaciones de diferencias finitas.
		Esquemas hacia adelante, hacia atrás, Crank-Nicolson para una
		ecuación de difusión no homogénea.~\cite[pág.~253]{Atkinson2009}.
	}                                                                                               \\
	\hline
	2                                                                                             &
	\cellteoria{
		Solución de un problema de valor inicial.
		Buen planteamiento de la solución.
		Solución generalizada de un problema de valor inicial.
		Consistencia.
		Convergencia.
		Estabilidad.~\cite[pág.~260]{Atkinson2009}.
	}                                                                                             &
	\celllab{
		Derivación de esquemas de diferencias para la ecuación de
		difusión, onda, transporte~\cite{Shiach2020,Strikwerda2004,Mortensen2023}.
		Simulaciones de la ecuación de transporte~\cite{Aznaran2025}.
	}                                                                                             &
	\cellteoria{
		Teorema de equivalencia de Lax.
		Orden de convergencia.
		Esquema de dos niveles.
		Revisión de consistencia, convergencia y estabilidad.
		Estimación de error~\cite[pág.~266]{Atkinson2009}.
	}                                                                                               \\
	\hline
	3                                                                                             &
	\cellteoria{
		Métodos de diferencias finitas para problemas elípticos.
		Funciones malla y Operadores de diferencia finita.
		Propiedades de operadores de diferencia.
		Punto malla interior, frontera, clausura~\cite[pág.~664]{Salgado2022}.
	}                                                                                             &
	\cellevaluacion{Práctica Calificada $1$}                                                      &
	\cellteoria{
		Propiedad de aproximación.
		Producto interno discreto.
		Operador de consistencia.
		Convergencia.
		Teorema de Lax.
		Problema de Poisson 1D.
		Principio del máximo discreto.
		Problema de Poisson 2D.~\cite[pág.~671]{Salgado2022}
	}                                                                                               \\
	\hline
	4                                                                                             &
	\cellteoria{
	Método de Galerkín.
	Lema de Cea.
	Teorema de Lax Milgram.
	Base de Lagrange. {[}2{]}
	}                                                                                             &
	\celllab{
	Esquema de 5 puntos para Poisson 1D y esquema de 9 puntos para Poisson 2D.
	Solución por métodos iterativos. {[}3{]}
	}                                                                                             &
	\cellteoria{
	Construcción de la matriz de rigidez y del vector de carga en el método de Galerkin. {[}2{]}
	}                                                                                               \\
	\hline
	5                                                                                             &
	\cellteoria{
	Motivación.
	Ecuación de transporte.
	Ondas acústicas.
	Flujo de tráfico.
	Aguas poco profundas.
	Ecuaciones de Euler. {[}4{]}
	}                                                                                             &
	\cellevaluacion{Práctica Calificada $2$}                                                      &
	\cellteoria{
	Leyes de conservación hiperbólicas escalar y sistemas.
	Solución clásica, débil y entrópica.
	Ondas de choque. {[}4{]}
	}                                                                                               \\
	\hline
	6                                                                                             &
	\cellteoria{
	Condición CFL.
	Ecuación de transporte.
	Ecuación de Burgers.
	Consistencia.
	Estabilidad von Neumann.
	Esquemas Upwind,
	FTCS,
	Lax-Friedrichs,
	Leapfrog.
	Problema de Riemann.{[}5{]}}                                                                  &
	\celllab{
	Ejercicios sugeridos:
	3.1, 3.2, 3.3, 3.4, 3.5, 3.6, 3.7, 3.8. {[}6{]}
	}                                                                                             &
	\cellteoria{
	Esquema BTCS, Lax-Wendroff, Crank-Nicolson, Du Fort-Frankel. {[}5{]}
	}                                                                                               \\
	\hline
	7                                                                                             &
	\cellteoria{
	Ecuaciones de Agua Poco Profundas 1D.
	Problema de Dam Break.
	Esquema de Mac-Cormack. {[}3{]}
	}                                                                                             &
	\celllab{
	Ejercicios:
	Detallar ejemplos sección 5.1 {[}7{]} Ejercicios capítulo 29 {[}2{]}
	}                                                                                             &
	\cellteoria{Presentación del Proyecto. Parte 1}                                                 \\
	\hline
	8                                                                                             &
	\multicolumn{3}{c|}{\cellevaluacion{Examen Parcial}}                                            \\
	\hline
	9                                                                                             &
	\cellteoria{Propiedad monótona de esquemas. Teorema de Godunov. {[}6{]}}                      &
	\celllab{Tutorial de Clawpack}                                                                &
	\cellteoria{Método de Volúmenes Finitos. Problema de Dirichlet. Estimación de error.}           \\
	\hline
	10                                                                                            &
	\cellteoria{Problemas parabólicos e hiperbólicos. {[}4{]}}                                    &
	\cellevaluacion{Práctica Calificada $3$}                                                      &
	\cellteoria{Esquema numérico para caso lineal y no lineal. {[}4{]}}                             \\
	\hline
	11                                                                                            &
	\cellteoria{Método de Godunov y variantes. {[}4{]}}                                           &
	\celllab{Ejercicios sugeridos}                                                                &
	\cellteoria{Teorema de Lax Wendroff. {[}6{]}}                                                   \\
	\hline
	12                                                                                            &
	\cellteoria{Resolvente de Riemann aproximado. {[}4{]}}                                        &
	\cellevaluacion{Práctica Calificada $4$}                                                      &
	\cellteoria{Método de direcciones alternadas. {[}6{]}}                                          \\
	\hline
	13                                                                                            &
	\cellteoria{Método WENO y ENO. {[}6{]}}                                                       &
	\celllab{Ejercicios sugeridos}                                                                &
	\cellteoria{Técnicas de descomposición de flujo}                                                \\
	\hline
	14                                                                                            &
	\cellteoria{Esquema de variación total decreciente (TVD) y consistente con entropía. {[}6{]}} &
	\celllab{Ejercicios sugeridos}                                                                &
	\cellteoria{Métodos de volúmenes finitos en mallas no estructuradas}                            \\
	\hline
	15                                                                                            &
	\cellteoria{Aplicaciones. {[}4{]}}                                                            &
	\cellevaluacion{Práctica Calificada $5$}                                                      &
	\cellteoria{Presentación del Proyecto. Parte II}                                                \\
	\hline
	16                                                                                            &
	\multicolumn{3}{c|}{\cellevaluacion{Examen Final}}                                              \\
	\hline
	17                                                                                            &
	\multicolumn{3}{c|}{\cellevaluacion{Examen Sustitutorio}}                                       \\
	\hline
\end{longtable}

\printbibliography[title={\color{Primary}Referencias}]

\end{document}
